% Issue #30 (Documentar a análise léxica), concluída; migrada do docs/analise-lexica.md.
% Issue #29: revisar contra o lexer.l final (#17 e #19 podem acrescentar tokens).

\subsection{Léxico}
\label{sec:lexico}

O analisador léxico de Lastro fica em \cod{src/lexer/lexer.l}, e o flex o
transforma na função \cod{yylex()} \cite{levine2009}. Quem comanda a leitura é o
analisador sintático, que pede um token por vez. A cada pedido, o léxico lê a
entrada padrão até fechar o próximo lexema e devolve o código do token. Se o
token tem valor, como um identificador ou um número, o lexema vai para
\cod{yylval.str}. A opção \cod{\%option yylineno} guarda a linha corrente, e
todas as mensagens de erro do compilador usam esse número, sejam elas léxicas,
sintáticas ou semânticas. O léxico só consulta a tabela de tipos, para devolver
\cod{TYPE\_ID} quando o identificador é o nome de um tipo. Ele entrega só o texto
do identificador, e o analisador sintático decide se aquilo é uma declaração ou
um uso, como fixamos no Problema~6.

\subsubsection*{Tokens e expressões regulares}

A especificação parte de três definições auxiliares: \cod{digito = [0-9]},
\cod{letra = [a-zA-Z]} e \cod{ident = \{letra\}(\{letra\}|\{digito\}|\_)*}.
A Tabela~\ref{tab:tokens} lista os tokens na mesma ordem em que as regras
aparecem no arquivo.

\begin{table}[ht]
\centering
\caption{Tokens de Lastro, na ordem das regras do \cod{lexer.l}}
\label{tab:tokens}
\small
\begin{tabularx}{\textwidth}{@{}l>{\raggedright\arraybackslash}X>{\raggedright\arraybackslash}X@{}}
\toprule
Classe & Lexema ou expressão regular & Token \\
\midrule
Palavra reservada & \cod{const proc func} & \cod{CONST PROC FUNC} \\
Palavra reservada & \cod{int float bool} & \cod{INT FLOAT BOOL} \\
Palavra reservada & \cod{if else while for to} & \cod{IF ELSE WHILE FOR TO} \\
Palavra reservada & \cod{switch case default} & \cod{SWITCH CASE DEFAULT} \\
Palavra reservada & \cod{break continue return ref} & \cod{BREAK CONTINUE RETURN REF} \\
Literal lógico & \cod{true false} & \cod{TRUE FALSE} \\
Atribuição & \cod{:=} & \cod{ASSIGN} \\
Relacional & \cod{= != <= < >= >} & \cod{EQ NEQ LE LT GE GT} \\
Aritmético & \cod{+ - * /} & \cod{PLUS MINUS STAR SLASH} \\
Lógico & \cod{\&\& || !} & \cod{AND OR NOT} \\
Tipo de retorno & \cod{->} & \cod{ARROW} \\
Delimitador & \cod{( ) [ ] \{ \}} & \cod{LPAREN RPAREN LBRACKET RBRACKET LBRACE RBRACE} \\
Separador & \cod{, ; .} & \cod{COMMA SEMI DOT} \\
Literal de texto & \cod{\textbackslash"[\textasciicircum\textbackslash"\textbackslash\textbackslash\textbackslash n]*\textbackslash"} & \cod{LITERAL\_TEXTO} \\
Texto sem fechar & \cod{\textbackslash"[\textasciicircum\textbackslash"\textbackslash\textbackslash\textbackslash n]*} & \cod{ERRO} \\
Literal real & \cod{\{digito\}+"."\{digito\}+} & \cod{LITERAL\_FLOAT} \\
Literal inteiro & \cod{\{digito\}+} & \cod{LITERAL\_INT} \\
Identificador & \cod{\{ident\}} & \cod{ID} ou \cod{TYPE\_ID} \\
Erro & \cod{.} (qualquer outro caractere) & \cod{ERRO} \\
\bottomrule
\end{tabularx}
\end{table}

A linguagem tem 20 palavras reservadas, todas em minúsculas. Como ela diferencia
maiúsculas de minúsculas, \cod{If} e \cod{WHILE} são identificadores comuns.
\cod{ler} e \cod{escrever} também são identificadores: chegam ao analisador
sintático como \cod{ID} e entram na tabela de símbolos como subprogramas
embutidos, de modo que a análise semântica confere uma chamada a eles do mesmo
jeito que confere qualquer outra.

Todo identificador começa com uma letra ASCII e continua com letras, dígitos ou
sublinhado. Deixamos as letras acentuadas de fora porque o código de sistemas
financeiros conversa com bibliotecas C e com sistemas legados (Problema~1), e um
nome como \cod{preço} não teria equivalente direto no C gerado. O literal
inteiro é uma sequência de dígitos, sem sinal. Em \cod{-5}, o \cod{-} é o
operador unário da gramática. O literal real exige dígitos dos dois lados do
ponto (\cod{2.5}; \cod{2.} e \cod{.5} não são aceitos). Em \cod{2.5}, as regras
de inteiro e de real casam ao mesmo tempo, e o casamento mais longo escolhe
\cod{LITERAL\_FLOAT}. Valores monetários continuam sendo guardados como inteiros
(em centavos ou pontos-base), como decidido no Problema~1. O \cod{float} existe
para os cálculos não monetários e para a coerção de \cod{int} para \cod{float}
exigida na avaliação.

O literal de texto fica entre aspas, na mesma linha e sem \cod{\textbackslash},
e vai para o C exatamente como foi escrito. A linguagem não tem tipo
\cod{string}: o texto só é aceito como argumento de \cod{escrever}, para as
mensagens exigidas nos problemas de avaliação. Quando as aspas não fecham na
linha, ou quando aparece uma \cod{\textbackslash}, só a regra de texto sem
fechar casa, e o léxico reporta \cod{texto sem aspas de fechamento}.

O léxico já reconhece \cod{SWITCH}, \cod{CASE}, \cod{DEFAULT} e \cod{DOT}, mas a
gramática ainda não usa esses tokens em nenhuma produção. Eles ficam reservados
para o comando \cod{switch} e para o acesso a campo de registro, e o yacc os
aponta como terminais não usados. Quando um deles aparece num programa, o léxico
o aceita e o analisador sintático acusa erro de sintaxe.

\subsubsection*{Casamento mais longo e palavras reservadas antes de identificador}

O flex desempata com duas regras \cite{aho2007}. Pela primeira, o casamento mais
longo (\emph{longest match}), vence a expressão que consome mais caracteres a
partir da posição atual. Se duas expressões casam o mesmo número de caracteres,
vale a segunda regra: ganha a que foi declarada primeiro no arquivo.

Com o casamento mais longo, os operadores que começam igual não precisam de
tratamento especial. Em \cod{<=}, tanto \cod{"<="} quanto \cod{"<"} casam, mas a
primeira leva dois caracteres e ganha. O mesmo vale para \cod{:=} e \cod{=},
\cod{>=} e \cod{>}, \cod{!=} e \cod{!}, \cod{->} e \cod{-}, e por isso a ordem
dessas regras no arquivo não muda nada. O comentário segue a mesma lógica:
\cod{"//".*} sempre consome mais que \cod{"/"}, então duas barras seguidas nunca
viram duas divisões.

Nas palavras reservadas, quem decide é a ordem. \cod{while} casa \cod{"while"} e
também \cod{\{ident\}}, com o mesmo comprimento, e só é lido como palavra
reservada porque essas regras estão antes da regra de identificador no arquivo.
Se a ordem fosse invertida, o léxico nunca devolveria uma palavra reservada. O
casamento mais longo cobre o caso inverso. Em \cod{ifx}, a regra \cod{"if"} casa
dois caracteres e \cod{\{ident\}} casa três, então \cod{ifx} é um identificador,
e o programador pode dar a uma variável um nome que começa com uma palavra
reservada.

\subsubsection*{Comentários, espaços e erro léxico}

O léxico reconhece comentários e espaços e os descarta sem gerar token. O
comentário vai de \cod{//} até o fim da linha (\cod{"//".*}), porque o \cod{.}
do flex não casa a quebra de linha. A linguagem não tem comentário de bloco,
então ninguém esquece um comentário aberto pelo caminho, e cada linha pode ser
lida sozinha, o que atende à prioridade de legibilidade do Problema~1. A
expressão \cod{[ \textbackslash t\textbackslash r\textbackslash n]+} cobre os
espaços, que só separam lexemas. O \cod{\textbackslash r} está ali para aceitar
arquivos com quebra de linha do Windows. As quebras de linha não geram token,
mas o \cod{yylineno} as conta.

A última regra, \cod{.}, pega qualquer caractere que nenhuma outra regra
aceitou: \cod{@}, \cod{\#}, um \cod{:} sozinho, um \cod{\&} ou \cod{|} sozinho,
ou cada byte de uma letra acentuada. A ação chama
\cod{yyerror("caractere inválido")}, que imprime a linha e o lexema, e devolve o
token \cod{ERRO}. Nenhuma produção usa esse token, então o analisador sintático
rejeita o programa no mesmo ponto e o compilador sai com código 1. Com
\cod{tests/invalido/05\_caractere\_invalido.txt}, a saída é:

\begin{lstlisting}[style=saida]
linha 3: caractere inválido perto de '@'
linha 3: syntax error perto de '@'
REJEITADO: programa contém erro
\end{lstlisting}

Dois dos erros previstos no Problema~6 formam sequências de tokens válidas e por
isso passam pelo léxico. Quem os detecta é o analisador sintático. O literal
fracionário incompleto \cod{10.} vira \cod{LITERAL\_INT DOT}, e a mensagem é
\cod{syntax error perto de '.'}. O literal mal formado \cod{3x} vira
\cod{LITERAL\_INT ID}, com a mensagem \cod{syntax error perto de 'x'}. Nos dois
casos o programa é rejeitado na linha certa.
